\begin{table}[!htbp]\small\setlength{\tabcolsep}{4pt}
\caption{Coût d'une baisse du prélèvement forfaitaire unique (Md€ de 2025, option non retenue)}\label{tab:t5}
\begin{tabularx}{\linewidth}{@{}Xrrrr@{}}
\toprule
\textbf{Option} & \textbf{Statique} & \textbf{Après réaction} & \textbf{Fourchette} & \textbf{Taux cumulé} \\
\midrule
Prélèvement forfaitaire unique à 30~\% & $-$1,5 & $-$1,1 & $-$1,4 à $-$1,0 & 47,5~\% \\
Prélèvement forfaitaire unique à 25~\% & $-$5,2 & $-$3,8 & $-$4,8 à $-$3,4 & 44~\% \\
\bottomrule
\end{tabularx}
\par\vspace{2pt}{\raggedright\scriptsize\color{gris}Source : DGFiP (dividendes, intérêts et plus-values déclarés au titre de 2024, 80 Md€) ; élasticité du revenu déclaré de 0,22 à 0,96 (Lefebvre, Lehmann et Sicsic) ; Sénat pour le rendement constaté de la hausse de 2026. Taux cumulé : impôt sur les sociétés à 25~\% puis prélèvement sur le dividende, à comparer aux 48~\% du haut du nouvel impôt sur le revenu.\par}
\end{table}
